\section{A Lagrangian immersion with controlled cohomology}
\label{sec:topology}

We begin the proof of Proposition~\ref{prop:realization} by constructing
a geometric representative of its input class $a$.  Two features of that representative will control its
Floer theory: ambient classes detect its degree-two cohomology, and its
self-intersections contribute only odd-degree generators.  We establish
these features together.  All homology and cohomology in this section have
rational coefficients unless otherwise indicated.

\begin{proposition}\label{prop:immersion}
Let $D=Y\times T_g$, $\omega$, and $n$ be as in Section~\ref{sec:setup}, with
$g$ even and $n\geq12$.  Suppose that $a\in H^n(D,\Q)$ satisfies
$\omega a=0$, and that its annihilator
\[
 \mathcal Z=\{x\in H^2(D,\Q):xa=0\}
\]
satisfies the injectivity hypothesis \eqref{setup:wedge-injective}.
There exist a closed connected spin $n$-manifold $L$, a positive integer
$\nu$, and a Maslov-zero Lagrangian immersion $i:L\to D$ such that
\begin{equation}\label{topo:class}
 \PD\,i_*[L]=\nu a,
\end{equation}
\begin{equation}\label{topo:h2}
 i^*:H^2(D,\Q)/\mathcal Z\xrightarrow{\ \sim\ }H^2(L,\Q),
\end{equation}
and
\begin{equation}\label{topo:push}
 i_*:H_{n-2}(L,\Q)\longrightarrow H_{n-2}(D,\Q)
 \quad\text{is injective}.
\end{equation}
The immersion is surjective on fundamental groups and has only transverse
double points, with no triple points.  With an integer grading whose
orientation agrees with the orientation of $L$, both ordered branch
indices at every double point are odd.
\end{proposition}

We first realize the class with the required tangential structure and low
cohomology.  The immersion theorem then supplies a graded Lagrangian
representative.  To control the double-point signs, we place all of its
crossings inside a contractible exact domain and apply the flexibility
results of Eliashberg--Murphy and Ekholm--Eliashberg--Murphy--Smith.  This
adapts the local-domain construction of \cite[Section~3]{A8}; the choice
of filling classes below also treats the nonzero $\pi_2$ of the K3 factor.

\subsection{Tangential realization and low-dimensional surgery}

We seek a representative whose degree-two cohomology is exactly
$H^2(D,\Q)/\mathcal Z$ and whose complexified tangent bundle is
identified with the ambient tangent bundle.  We impose these two
conditions by lifting the representing map: a torus bundle kills
the classes in $\mathcal Z$, and a spin pullback records the required
stable tangent identification.  The following calculation shows that
the lift introduces no further degree-two classes; surgery will then
transfer this degree-two description to the representing manifold.

Choose independent integral classes $z_1,\ldots,z_r$ spanning
$\mathcal Z$ over $\Q$, and let
$\pi_0:\mathscr T_0\to D$ be the principal $r$-torus bundle with these
Chern classes.  Choose the natural stable $\SU$ structure on $TD$, and
form the homotopy pullback
\[
 \mathscr T=\mathscr T_0\times^{h}_{B\SU}B\Spin,
 \qquad \pi:\mathscr T\longrightarrow D,
\]
where $B\Spin\to B\SU$ is complexification.  The tautological stable
real bundle $\xi$ on $\mathscr T$ has virtual rank $n$, is spin, and
comes with a stable identification
\begin{equation}\label{topo:tautological}
 \xi\otimes\C\simeq\pi^*TD
\end{equation}
that retains the determinant trivializations.

\begin{lemma}\label{topo:lift}
Pullback induces
\[
 H^2(\mathscr T,\Q)\simeq H^2(D,\Q)/\mathcal Z.
\]
Moreover, $\PD^{-1}a\in H_n(D,\Q)$ is in the image of
$\pi_*:H_n(\mathscr T,\Q)\to H_n(D,\Q)$.
\end{lemma}

\begin{proof}
The compact K\"ahler manifold $D$ is formal over $\R$
\cite{DeligneGriffithsMorganSullivan}.  We first compute over $\R$.
A real model for its torus bundle is obtained from $H^*(D,\R)$
by adjoining degree-one
generators $u_1,\ldots,u_r$ with $du_j=z_j$.  In degree two, the
horizontal classes give $H^2(D,\R)/(\mathcal Z\otimes\R)$.
The differential on
$H^1(D,\R)\otimes\langle u_1,\ldots,u_r\rangle$ is the scalar
extension of the map in \eqref{setup:wedge-injective}, and is
injective.  The differential
on $\bigwedge^2\langle u_1,\ldots,u_r\rangle$ is also injective:
\[
 d(u_j u_l)=z_j u_l-z_l u_j,
\]
and the $z_j$ are independent.  The images of these two differentials
have different vertical degrees, so cannot cancel.

The fiber of $B\Spin\to B\SU$ is two-connected.  Over $\R$,
complexification takes each even-index Chern generator to a nonzero
multiple of a free Pontryagin generator, and every odd-index Chern
generator to zero.  In the relative Koszul model for the second
pullback, eliminate the even-index pairs.  What remains is the torus
bundle model with generators of degrees $4j+1$, $j\geq1$, whose
differentials are the odd Chern classes of $TD$.  These classes vanish:
$TD$ is the sum of the tangent bundle of a K3 surface, when present,
and a trivial bundle.  In particular the second pullback changes
neither the preceding degree-two calculation nor its pullback map.

Define a functional on the degree-$n$ part of this model by pairing
a horizontal class $v\in H^n(D,\R)$ with $a$, and setting the
functional equal to zero on terms with positive vertical degree.
It is a cycle in the dual complex.  Indeed, a differential with a
horizontal term is a sum of multiples of the $z_j$, and
$a z_j=0$; the remaining new generators have zero differential.
Its induced real homology class maps to $\PD^{-1}a$.
The cohomology pullback and homology pushforward are defined over
$\Q$.  Faithfulness of scalar extension from $\Q$ to $\R$ therefore
descends the asserted kernel and isomorphism, as well as membership
of the rational class $\PD^{-1}a$ in the image of pushforward.
\end{proof}

\begin{lemma}\label{topo:source}
There are a closed connected spin $n$-manifold $L_0$, a map
$f:L_0\to D$, and a positive integer $\nu$ satisfying
\eqref{topo:class}--\eqref{topo:push} with $f$ in place of $i$.
They admit a stable identification
\begin{equation}\label{topo:stable}
 TL_0\otimes\C\simeq f^*TD
\end{equation}
with the specified determinant homotopy.  The map is surjective on
fundamental groups, and $L_0$ has two additional free circle generators
mapping trivially to $D$.  After fixing $\nu$, its Euler characteristic
can be made arbitrarily large while preserving these assertions.
\end{lemma}

\begin{proof}
Use the spin bundle $\xi$ on $\mathscr T$, including
\eqref{topo:tautological}, as a stable tangential structure.
With our virtual-rank-$n$ convention, the group of $n$-dimensional
bordism classes carrying this structure is
$\pi_0\operatorname{Th}(-\xi)$ by the Pontryagin--Thom construction;
see \cite{Thom1954} and \cite[Section~2.3]{RathelFournier2024}.
Rational stable homotopy and the Thom isomorphism give
\[
 \pi_0\operatorname{Th}(-\xi)\otimes\Q
 \simeq H_0(\operatorname{Th}(-\xi);\Q)
 \simeq H_n(\mathscr T;\Q).
\]
Thus a positive integral multiple of the lift in Lemma~\ref{topo:lift}
is represented
by a closed manifold mapping to $\mathscr T$, with its stable tangent
bundle identified with the pullback of $\xi$.  Composing with $\pi$
gives \eqref{topo:stable} and the required fundamental class.  We now
modify the representative by bordisms over $\mathscr T$.

Put $\Gamma=\pi_1(\mathscr T)$.  The torus-bundle homotopy sequence
exhibits its fundamental group as an extension of $\pi_1(D)$ by a
quotient of $\Z^r$; the second pullback has two-connected fiber.
Consequently $\Gamma$ is polycyclic, and $\Gamma\to\pi_1(D)$ is
surjective.  The space $\mathscr T$ has a CW model with finite skeleta
in the degrees needed here: the base and torus bundle are compact,
and the characteristic-class pullback uses classifying spaces and
loop spaces of finite type.  The group ring $\Z\Gamma$ is
Noetherian \cite{Hall1954}.

Join the components of the representative by one-handles and add
loops mapping to generators of $\Gamma$.  Its fundamental-group map
is then surjective.  Its kernel is normally finitely generated,
since both groups are finitely presented.  Surgery on disjoint
embedded representatives of normal generators makes that map an
isomorphism.  Each surgery uses a null homotopy over $\mathscr T$;
a twist of the normal framing on a circle makes the stable tangential
data extend.

Next make the map on $\pi_2$ surjective by adding
$S^2\times S^{n-2}$ summands mapped through $S^2$, for a finite set
of module generators of the missing classes.  Such a pulled-back
spin bundle is stably trivial over $S^2$, so these additions carry
the prescribed tangential data.  The kernel on $\pi_2$ is finitely
generated as a $\Z\Gamma$-module: this follows from the finite
cellular chain modules of the universal covers and Noetherianity.
Represent its generators by disjoint embedded two-spheres and
perform surgery on them, using their null homotopies over
$\mathscr T$.  Their normal bundles are trivial: the spin condition
makes their stable normal bundles trivial, and the rank is in the
stable range.  The stable framing extends over surgery because
$\pi_2\SO=0$.  Since $n\geq12$, the embedding dimensions and the
reverse-handle indices preserve the already established
fundamental-group assertion.  These operations make $\pi_2$ an
isomorphism as well.  The resulting map to $\mathscr T$ consequently
induces an isomorphism on $H^2$.

Now take connected sums with two copies of $S^1\times S^{n-1}$,
and with as many copies of $S^4\times S^{n-4}$ as desired, mapping
these summands constantly to $\mathscr T$.  These products are
stably parallelizable and preserve the tangential structure and
pushed fundamental class.  They do not change $H^2$.  The first
two summands provide the extra null-image circles; each subsequent
summand raises the Euler characteristic by two.  Lemma~\ref{topo:lift}
therefore proves \eqref{topo:h2}.

To verify \eqref{topo:push}, write any $u\in H_{n-2}(L_0,\Q)$ as
$f^*x\cap[L_0]$ using Poincar\'e duality and the surjectivity in
\eqref{topo:h2}.  The projection formula gives
\[
 f_*u=x\cap f_*[L_0]=\nu\,\PD^{-1}(xa).
\]
If this is zero, then $x\in\mathcal Z$ and $f^*x=0$, so $u=0$.
\end{proof}

Since $\omega\in\mathcal Z$, the map just constructed satisfies
$f^*[\omega]=0$.  Destabilize \eqref{topo:stable} in the unitary
stable range over the $n$-dimensional manifold $L_0$, retaining the
determinant homotopy.  Restriction to the real tangent bundle gives
a formal Lagrangian immersion of Maslov class zero.  The
Gromov--Lees theorem realizes it by a Lagrangian immersion homotopic
to $f$ \cite{GromovLees,RathelFournier2024}.  We continue to call
this immersion $f$, and choose its integer grading with the
orientation of $L_0$.

\subsection{Chosen fillings of zero area}

The remaining task is to alter the double points without changing
the cohomological properties.  To do this inside a contractible
exact domain, the loops formed by the two branches of each crossing
must admit chosen fillings of zero symplectic area.  No uniqueness
of a filling area is asserted or required.

\begin{lemma}\label{topo:zero-actions}
After Lagrangian perturbations, $f$ is self-transverse without triple
points, and there are pairwise disjoint embedded arcs in $L_0$
joining the two preimages of each double point such that every
image loop has a chosen filling in $D$ of symplectic area zero.
The perturbations preserve the homotopy class and integer grading.
\end{lemma}

\begin{proof}
Choose fillings for the two extra null-image circles from
Lemma~\ref{topo:source}.  In an immersed cotangent neighborhood,
replacing the zero section by a small closed one-form changes each
chosen filling area by the corresponding period, up to a fixed
sign.  The two periods can be adjusted independently.  Choose
the resulting areas to generate a dense subgroup of $\R$.
Concatenations and reversals of these circles are given the
concatenated and reversed fillings, so their chosen areas add.
Small exact perturbations now make the immersion self-transverse
and remove triple points, without changing these periods.

For a double point $f(a_j)=f(b_j)$, begin with a path from $a_j$
to $b_j$.  Surjectivity on fundamental groups allows insertion of
a source loop that makes its image contractible.  Choose a filling.
By inserting combinations of the two null-image circles with their
chosen fillings, make its area arbitrarily small.  Choose disjoint
source patches at the $a_j$, avoiding every other double-point
preimage.  An exact perturbation by $dF_j$, where $F_j$ is constant
near $a_j$ and zero off its patch, changes this path's filling area
by $F_j(b_j)-F_j(a_j)$, up to the same sign convention.  It fixes
the sheet near the crossing.  Thus it makes the area exactly zero.
The endpoint formula remains valid if the path visits the patch
several times.  With the patches fixed, a uniform sufficiently
small bound on these perturbations preserves all crossings and
creates no new ones; take the preceding residual areas within that
bound.

Finally, homotopies of the source paths relative to their endpoints
preserve the chosen areas, because $f^*\omega=0$ as a differential
form.  General position makes the finitely many paths embedded,
mutually disjoint, and disjoint from the other double-point
preimages.  Grading and spin data transport through the
Lagrangian perturbations.
\end{proof}

\subsection{A contractible exact neighborhood}

\begin{lemma}\label{topo:exact-domain}
There is a compact contractible Weinstein domain
$(P,\lambda)\subset(D,\omega)$, with $d\lambda=\omega|_P$, such that
$W=f^{-1}(P)$ is a disk
containing all double-point preimages.  Its primitive satisfies
$f^*\lambda|_W=0$, and its Liouville field is tangent to the
immersed sheets.  The immersion is cylindrical along $\partial W$.
At one interior point, $W$ has a single linear sheet in a radial
Darboux ball.
\end{lemma}

\begin{proof}
The source and target neighborhoods have different topology.  A tree
joining the double-point preimages has a disk as its source thickening,
but its image acquires a loop at each double point.  We first thicken
that image to a Weinstein zero- and one-handlebody $P_1$.  We then
attach two-handles along its free generators, with cores disjoint from
the immersion.  The target becomes contractible while its preimage
remains a disk.  We give the primitive, collar, and filling details
needed for this construction, adapting \cite[Lemma~3.5]{A8}.


\emph{Thickening the tree and its image.}
Extend the arcs of Lemma~\ref{topo:zero-actions} to an embedded tree
$\Theta\subset L_0$ by connector arcs avoiding all other sheet
intersections.  Its image $G=f(\Theta)$ has precisely the
identifications at the double-point pairs.  Each chosen source arc
becomes one of the free graph cycles, with its previously chosen
zero-area filling.  Subdivide an ordinary edge to obtain a vertex
with only one preimage.
Figure~\ref{fig:topo-spines} illustrates the resulting spines.

\begin{figure}[htbp]
\centering
\setlength{\unitlength}{1pt}
\begin{picture}(420,105)
 \put(10,53){\line(1,0){90}}
 \put(10,53){\circle*{4}}
 \put(35,53){\circle*{4}}
 \put(75,53){\circle*{4}}
 \put(100,53){\circle*{4}}
 \put(10,67){\makebox(0,0){$a_1$}}
 \put(35,67){\makebox(0,0){$b_1$}}
 \put(75,67){\makebox(0,0){$a_2$}}
 \put(100,67){\makebox(0,0){$b_2$}}
 \put(55,19){\makebox(0,0){source tree $\Theta$}}
 \put(55,5){\makebox(0,0){thickening: $W$}}
 \put(119,53){\vector(1,0){31}}
 \put(134,68){\makebox(0,0){$f$}}
 \put(177,53){\circle{32}}
 \put(241,53){\circle{32}}
 \put(193,53){\circle*{4}}
 \put(225,53){\circle*{4}}
 \put(193,53){\line(1,0){32}}
 \put(209,19){\makebox(0,0){image graph $G$}}
 \put(209,5){\makebox(0,0){thickening: $P_1$}}
 \put(271,53){\vector(1,0){31}}
 \put(321,53){\color{black!15}\scalebox{2.1333}{\circle*{15}}}
 \put(385,53){\color{black!15}\scalebox{2.1333}{\circle*{15}}}
 \put(321,53){\circle{32}}
 \put(385,53){\circle{32}}
 \put(337,53){\circle*{4}}
 \put(369,53){\circle*{4}}
 \put(337,53){\line(1,0){32}}
 \put(321,53){\makebox(0,0){$e^2_1$}}
 \put(385,53){\makebox(0,0){$e^2_2$}}
 \put(353,19){\makebox(0,0){filled spine}}
 \put(353,5){\makebox(0,0){thickening: $P$}}
\end{picture}
\caption{Schematic spines for two double points.  Identifying
$a_j$ with $b_j$ turns a contractible source tree into a graph with
loops.  Two-handles kill these loops in the target; their filling
interiors miss the immersion, so the preimage remains the disk $W$.
The picture represents the abstract spines, not an embedding in $D$.}
\label{fig:topo-spines}
\end{figure}

At each vertex take a small radial Darboux ball.  The single sheet,
or the two transverse sheets at a double point, can be straightened
simultaneously.  For two sheets, first make one horizontal; the other
is a gradient graph over the vertical plane, and a symplectic shear
straightens it while fixing the first.  Arrange the incident tree
edges as distinct radial rays.  Thin Weinstein one-handles along the
remaining image arcs join these balls and carry strips in the
corresponding source sheets.

Here are compatible models for these attachments and the two-handles
used below.  In symplectic coordinates $(u,v,z)$, with
$u,v\in\R^j$ and $z\in\R^{2(n-j)}$, use the Liouville field
\begin{equation}\label{topo:handle-field}
 Z_j=(-u,2v,z/2),\qquad j=1,2.
\end{equation}
A box $|u|\leq1$, $|(v,z)|\leq\varepsilon$ has negative face near
$|u|=1$, isotropic core $v=z=0$, and a quadratic Lyapunov function
with one critical point of index $j$.  For $j=1$, the subspace
$v=0$ with $z$ in a fixed Lagrangian plane is an invariant
Lagrangian strip.

The attaching contact identifications can be chosen strict for the
specified contact forms.  Indeed, near an isotropic attaching locus
choose a slice tangent to the contact plane along that locus.
The Reeb flow writes the contact form as $dt+\beta$ on a
neighborhood of the slice.  Match the slices by the symplectic
isotropic neighborhood theorem with the prescribed tangent and
symplectic normal data.  The resulting difference of primitives is
a closed form that restricts to zero on the attaching locus.
It is $dh$, with $h=0$ there, since the neighborhood retracts onto
the locus.  Adjusting the Reeb coordinate by $h$ makes the contact
identification strict.  The same argument applies to the Legendrian
patches in the sheets.  Extending by Liouville flow gives collars
preserving the primitives.

For the one-handles, choose transverse strip framings along the
source arcs and extend the prescribed end germs over each strip.
The relative Lagrangian neighborhood theorem extends these maps
symplectically with the collar germs fixed: after fixing tangent
data, the remaining linear lifts are contractible, and the tubular
construction and Moser argument are relative to the collars.
Shrinking the handles keeps them embedded, mutually disjoint away
from their attachments, and disjoint from every other sheet.
The model primitives glue through the strict collars.  After
rounding the positive corners, we obtain a Weinstein zero- and
one-handlebody $P_1$ with spine $G$.  Its preimage is a thickening
of the source tree, hence a disk, and its primitive vanishes on the
sheets.  The singled-out radial ball can be retained inside $P_1$.

\emph{Isotropic fillings outside the one-handlebody.}
The loops of $P_1$ now have to be killed without meeting any source
sheet.  The fillings must retain their chosen zero-area classes:
contractibility of the attaching circles alone would not control area
in the presence of $\pi_2(D)$.  Choose a free basis of $\pi_1(P_1)$
represented by disjoint isotropic circles in $\partial P_1$, avoiding the Legendrian boundary of the
source disk.  These representatives exist by the relative subcritical
contact isotropic embedding theorem
\cite[Theorem~12.4.1]{EliashbergMishachev2002}; formal tangent lines
can be moved into the contact distribution, and all the avoidance
conditions are in the general-position range.

Each attaching circle inherits a zero-area filling in $D$: join it
in $P_1$ to its graph cycle by an annulus and append the chosen
filling.  Both boundary primitive integrals vanish, so the annulus
has area zero.  We may represent this same filling class outside $P_1$.
In fact, recovering a $2n$-dimensional zero- and one-handlebody from
its boundary uses only handles of indices $2n$ and $2n-1$.
The inclusion of the exterior into $D$ therefore has no relative
homotopy in the dimensions of this disk and its deformation.
This preserves the chosen relative disk class, not just the
contractibility of the attaching circle.  Consequently its
symplectic area remains zero even when $\pi_2(D)\neq0$.
Prescribe the initial collar to be the outward Liouville cylinder.
General position relative to these collars gives embedded filling
disks with disjoint interiors, all missing $f(L_0)$.

We next make the disks isotropic while retaining those collars.
Use a prequantization circle bundle for the integral form $\omega$.
The collars lift horizontally.  In a trivialization over a filling
disk, the zero area gives winding number zero for its boundary
lift, so this lift extends over the disk.  The space of isotropic
ordered two-frames in the contact distribution retracts to
$\mathrm U(n)/\mathrm U(n-2)$ and is simply connected.  Thus the
prescribed isotropic derivative on the collar extends as a formal
isotropic monomorphism over the disk.  The ordinary real two-frame
space in rank $2n+1$ is two-connected, so the derivative of the
lifted disk is homotopic to that monomorphism relative to the
collar.  The relative, $C^0$-dense subcritical contact embedding
$h$-principle realizes these data, fixing a slightly smaller collar.

Projection gives isotropic immersions of the disks into $D$.
They can be chosen embedded and mutually disjoint: contact
Hamiltonian variations near distinct lift points give arbitrary
independent projected variations, and $2+2<2n$.  The same argument
preserves avoidance of the immersed source because $2+n<2n$.
Smallness protects avoidance away from the fixed collars; the
collars themselves protect the attachments.

\emph{Attaching the two-handles.}
We now have embedded isotropic disks with disjoint interiors outside
$P_1$, prescribed outward collars, and no intersections with the
immersion.  Choose symplectic normal framings over these disks.  They
restrict to the contact symplectic normal framings needed at the attaching
circles.  Parameterize the attaching circles by parameterizations
extending over the disks.  The relative isotropic neighborhood
theorem and relative Moser argument now embed thin index-two
handles with the prescribed strict collars and model
\eqref{topo:handle-field}.  These handles miss the immersion.
Their primitives glue, and rounding corners produces $P$.

Its spine is the graph $G$ with a two-cell attached along each
free generator.  After collapsing a maximal tree, these attaching
maps form a free basis of a bouquet of circles; the resulting
complex is contractible.  Thus $P$ is contractible.  The two-handles
have not changed the preimage disk, its primitive-zero sheets,
or its outer cylindrical collar.  The radial single-sheet ball
also remains intact.
\end{proof}

We have now localized all double points in an exact simply connected
target.  This is what permits their algebraic signed count, rather
than a fundamental-group-valued intersection invariant, to control
the final replacement.

\subsection{Double-point signs and the integer grading}

\begin{proof}[Proof of Proposition~\ref{prop:immersion}]
Write $n=2k$.  For the immersion supplied by the preceding lemmas,
the self-intersection formula reads
\begin{equation}\label{topo:double-number}
 2I(f)=f_*[L_0]\cdot f_*[L_0]-(-1)^k\chi(L_0).
\end{equation}
Here the normal Euler number is $(-1)^k\chi(L_0)$, since the
normal bundle of an oriented Lagrangian is its tangent bundle
with this orientation factor.  In Lemma~\ref{topo:source}, choose
the Euler characteristic large enough that $I(f)$ is nonzero
and has sign $(-1)^{k-1}$.  Put $d=|I(f)|$.

Complete the domain $P$ of Lemma~\ref{topo:exact-domain} by its
positive cylindrical end, extend $W$ by its cylinder, and remove
the center of its single-sheet radial ball.  The punctured ball
is a negative Liouville end.  The source is now
$S^{n-1}\times\R$, with the standard real-plane Legendrian
sphere $\phi_0$ at that end.  The target is simply connected:
$P$ was contractible, and removing a point in dimension $2n$
does not change its fundamental group.

The stabilization traces of
\cite[Lemma~3.2 and proof of Lemma~3.4]{EEMS} have signs
\begin{equation}\label{topo:trace-signs}
 \phi\longrightarrow\phi_1:\ (-1)^{k-1},
 \qquad
 \phi_1\longrightarrow\phi:\ (-1)^k,
\end{equation}
where $\phi_1$ is a loose stabilization and the arrows run
outward along the symplectization.  Insert $d$ reverse traces
in separated collars near the puncture, from the $d$-fold
stabilization $\phi_d$ at the negative end to $\phi_0$ outward.
The new negative link is connected and loose, and its total
signed double count is zero by \eqref{topo:double-number}.

Eliashberg--Murphy's exact cancellation theorem
\cite[Theorem~2.3]{EliashbergMurphyCaps} now applies.  The target
is a simply connected Liouville manifold with a negative end;
the immersion is proper, cylindrical there, self-transverse,
and has finitely many double points.  It is exact because
$H^1(S^{n-1}\times\R;\R)=0$.  The loose component and vanishing
signed count were just arranged.  The theorem's extra width
condition concerns source dimension three, so is irrelevant
here.  We obtain a compactly supported Hamiltonian regular
homotopy to an embedding, leaving both cylindrical tails fixed.

In the unchanged negative tail insert $d$ forward traces from
$\phi_0$ to $\phi_d$, and fill the center by the real plane.
The resulting disk immersion has exactly $d$ double points,
all of sign $(-1)^{k-1}$.  Its noncylindrical part is compact.
A sufficiently long negative Liouville flow in the filled
positive completion brings that part inside $P$.  Backward
flow exists for all time; points of $P$ cannot leave it
backward, while the positive end is translated toward $P$.
The flow rescales the symplectic form and hence preserves
Lagrangianity.  The cylindrical set is unchanged, so its
intersection with the original outer collar is unchanged.
Trim at that collar and glue to the immersion outside $W$.

This replaces a source disk by a source disk with the same
collar.  A collar identification extends continuously over
the disks, and their maps into the contractible $P$ are
homotopic relative to the collar.  It follows that the pushed
fundamental class and the maps on degree-two and degree-$(n-2)$
(co)homology are unchanged, under the identifications through
the common complement.  The same complement carries every
first homology class and the fundamental-group map.
The boundary sphere has a unique spin structure, so spin extends
over the replacement disk.  The Maslov class remains zero.
Choose the extended integer grading with the
original orientation parity.

For graded oriented Lagrangian intersections in real source
dimension $2k$, a generator of degree $|x|$ has intersection
sign
\[
 (-1)^k(-1)^{|x|}.
\]
The sign $(-1)^{k-1}$ obtained above therefore makes every
branch degree odd.  Reversing the ordering replaces its degree
by $n-|x|$, which is also odd because $n$ is even.  The disk
replacement has only transverse double points and introduces
no triples or intersections with the unmodified source.
Together with Lemma~\ref{topo:source}, this proves all assertions.
\end{proof}
